\section{Exact Lean API map}

This appendix records the minimum interface needed to complete the formal
development.

\subsection{\code{PR} and terms}

The repository defines:
\begin{lstlisting}[language={}]
inductive PR : Nat -> Type
  | zero : PR 0
  | succ : PR 1
  | proj {k : Nat} (i : Fin k) : PR k
  | comp {m k : Nat} (f : PR m) (gs : Fin m -> PR k) : PR k
  | prec {k : Nat} (f : PR k) (g : PR (k + 2)) : PR (k + 1)

inductive PTm (n : Nat) : Type
  | var (i : Fin n) : PTm n
  | app {k : Nat} (f : PR k) (ts : Fin k -> PTm n) : PTm n
\end{lstlisting}
The recursion argument of \code{prec} comes first:
\[
\begin{aligned}
 \mathsf{prec}(f,g)(0,\vec x)&=f(\vec x),\\
 \mathsf{prec}(f,g)(S y,\vec x)
   &=g(y,\mathsf{prec}(f,g)(y,\vec x),\vec x).
\end{aligned}
\]

\subsection{\PRAAx}

The axiom family has constructors for projections, composition, zero and
successor recursion equations, \(Sx\ne0\), successor injectivity, and
quantifier-free induction:
\begin{lstlisting}[language={}]
PRAAx.proj
PRAAx.comp
PRAAx.precZero
PRAAx.precSucc
PRAAx.succNeZero
PRAAx.succInj
PRAAx.ind (phi : Fm PTm (n + 1)) (h : phi.QF)
\end{lstlisting}
The namespace \code{PRA} provides substitution-ready theorem wrappers:
\begin{lstlisting}[language={}]
PRA.proj
PRA.comp
PRA.precZero
PRA.precSucc
PRA.succNeZero
PRA.qfInd
\end{lstlisting}

The signature of \code{PRA.qfInd} is conceptually:
\[
\frac{
 \Gamma\vdash\varphi(0)
 \qquad
 \varphi(y),\Gamma^\uparrow\vdash\varphi(Sy)
 \qquad \varphi\text{ quantifier-free}
}{
 \Gamma\vdash\forall y\,\varphi(y)
}.
\]
The formalization should use this theorem directly rather than rebuild
induction from the raw axiom.

\subsection{Proof composition}

The deduction library supplies hypothesis, axiom, implication introduction
and elimination, conjunction operations, universal introduction and
elimination, and substitution.  Local node contracts should be packaged as
ordinary Lean values of the relevant \code{Prf} type.  The graph compiler is
then a meta-level recursive function that combines those values.

A useful first representation is:
\begin{lstlisting}[language={}]
structure LocalContract (n : Nat) where
  premise : Fm PTm n
  result  : Fm PTm n
  premiseQF : premise.QF
  resultQF  : result.QF
  proof : [premise] |-[(PRAAx)] result
\end{lstlisting}
The actual context notation and universe parameters should follow
\code{ConciseIPRA.Deduction}; the blueprint file keeps these details explicit.

\subsection{Native transfer}

The repository theorem
\begin{lstlisting}[language={}]
pra_iff_native :
  (Gamma |-[(PRAAx)] phi) <->
  NPrf (Gamma.map NFm.ofPRA) (NFm.ofPRA phi)
\end{lstlisting}
is the final bridge.  It should appear at the end of the dependency graph.
Using it earlier does not reduce the obligation to construct the left-hand
derivation.

\subsection{Recommended theorem names}

\begin{longtable}{p{.34\linewidth}p{.56\linewidth}}
\toprule
\textbf{Name} & \textbf{Statement}\\
\midrule
\endhead
\code{PRA.allBelow\_correct} &
Boolean fold equals \(1\) exactly when all bounded tests succeed.\\
\code{PRA.firstTrue\_correct} &
Bounded witness routine returns a valid least witness or the sentinel.\\
\code{PRA.tableRun\_zero}, \code{tableRun\_succ} &
Defining equations of the coded trace.\\
\code{PRA.traceReflection} &
Base/step/final local proofs imply the final trace verifier.\\
\code{PRA.reverseDemand\_adequate} &
Every upstream request dominates each outgoing demand.\\
\code{PRA.relationalPath} &
Finite viability yields a compatible path through a requested horizon.\\
\code{PRA.gapDecision} &
A fine approximation plus zero-or-gap evidence decides the alternative.\\
\code{PRA.contractDAGReflection} &
Local node proof objects compose to the global sink contract.\\
\code{PRA.cubicBisection} &
The cubic trace returns a regular residual name with quantitative uniqueness.\\
\bottomrule
\end{longtable}
